\documentclass[11pt,a4paper]{article}
\usepackage[utf8]{inputenc}
\usepackage[T1]{fontenc}
\usepackage[french]{babel}
\usepackage{amsmath,amssymb,amsthm,geometry,hyperref}
\geometry{margin=2.4cm}
\setlength{\emergencystretch}{3em}
\newtheorem{proposition}{Proposition}
\title{Supports auxiliaires gradués pour le transport FE\\
Construction discrète et portée du contre-exemple résolu}
\author{Étude GPU SLFE, fork keckj/rodin}
\date{16 novembre 2026}
\begin{document}
\maketitle

\section{Résolution locale et norme discrète}
Le contre-exemple initial transporte une bosse de largeur $0{,}0015$ sur
un maillage triangulaire fortement gradué. Le transfert FE vers un support
uniforme insuffisamment résolu supprime une partie de cette information.
Un support tensoriel gradué est considéré pour conserver la résolution
locale sans raffiner uniformément tout le domaine.

Le domaine est le tore unité. Deux partitions périodiques croissantes
définissent des rectangles de dimensions $h_i^x,h_j^y$. L'espace auxiliaire
$W$ est constitué de polynômes $Q_q$ discontinus. Les inconnues sont leurs
valeurs aux points de Gauss tensoriels, avec $q+1$ points par direction.
Si $\omega_a$ sont les poids de Gauss sur $[0,1]$, la norme est
\[
 \|w\|_W^2=\sum_{i,j,a,b} h_i^x h_j^y\omega_a\omega_b
                         |w_{ijab}|^2.
\]
La quadrature est exacte pour le carré d'un polynôme $Q_q$; cette norme
coïncide donc avec sa norme $L^2$. Les graduations testées sont fixes dans
le temps. La construction est un raffinement tensoriel, comportant des
bandes fines, et non une adaptation locale par quadtree.

\section{Transfert FE et translation sur chaque ligne}
Le transfert $R:V\to W$ est la projection $L^2$ des polynômes triangulaires
$P_p$ sur $Q_q$. Les intersections triangle--rectangle sont triangulées
et intégrées par une quadrature de degré au moins $p+2q$.
L'adjoint utilise les masses propres à chaque rectangle. Les noyaux
uniformes antérieurs sont conservés comme cas particulier et comparés
à leur version publiée.

Pour un cisaillement horizontal, chaque point transverse $y_{jb}$ définit
un déplacement $d_{jb}=\Delta t\,u(y_{jb})$. Sur cette ligne, les blocs du
transfert entre un intervalle source $I_s$ et un intervalle cible $I_i$ sont
\[
 B_{ia,sc}(d)=\frac{1}{h_i^x\omega_a}
 \int_{I_i\cap(I_s+d)}
 \ell_a\!\left(\frac{x-x_i}{h_i^x}\right)
 \ell_c\!\left(\frac{x-d-x_s}{h_s^x}\right)\,dx.
\]
L'intégrande a degré $2q$ et $q+1$ points de Gauss suffisent. La recherche
d'intervalles, la quadrature, les produits et les réductions sont effectués
sur GPU. Les intervalles sources et cibles peuvent appartenir à des
partitions différentes, à condition de conserver la même partition transverse.

\begin{proposition}
Pour tout déplacement réel, le transfert par lignes préserve les constantes
et la masse discrète et est contractant dans la norme tensorielle.
Son adjoint de masse est obtenu en échangeant les partitions source et cible
et en changeant le signe du déplacement.
\end{proposition}
\begin{proof}
Sur chaque ligne, le transfert est une translation unitaire suivie d'une
projection orthogonale dans l'espace polynomial cible. Les constantes sont
contenues dans les deux espaces. L'intégration contre la constante établit
la conservation de la masse sur la ligne. L'addition des inégalités de
contraction avec les poids positifs $h_j^y\omega_b$ donne la contraction
globale. Le changement de variable $x\mapsto x-d$ dans le produit scalaire
établit l'expression de l'adjoint.
\end{proof}

La dépendance en la coordonnée transverse est traitée par collocation aux
points de Gauss. La proposition concerne cet opérateur discret; elle
n'affirme pas l'exactitude de la projection bidimensionnelle du cisaillement
continu pour une vitesse transverse non polynomiale.

\section{Modes diagonaux du flot cellulaire}
La fonction de courant est $\psi=\sin^2(\pi x)\sin^2(\pi y)/\pi$.
Ses modes cartésiens sont des cisaillements directs. Pour les modes
diagonaux, les coordonnées sont $(\xi,\eta)=(x+y,y)$ et $(x-y,y)$.
Chaque support orienté est gradué autour de la position correspondante de
la bosse. Les transferts $C_+$ et $C_-$ réalisent le changement de
coordonnées par les projections sur lignes précédentes. Ils sont
contractants dans leurs normes discrètes respectives.

Un sous-pas diagonal sur le support principal est défini par
\[
 A_\pm=I+C_\pm^*(S_\pm-I)C_\pm.
\]
Il est contractant. En effet, pour une contraction $C$, l'application
$Jw=(Cw,(I-C^*C)^{1/2}w)$ est isométrique et
$I+C^*(S-I)C=J^*\operatorname{diag}(S,I)J$.
La racine carrée intervient uniquement dans cette preuve. Les propriétés
sur les constantes et la masse se transmettent aux sous-pas et à leur
composition. L'état principal $w$ est conservé entre les pas, avec retour
FE final $v=v_0+R^*(w-Rv_0)$.

Les sous-pas sont composés dans l'ordre symétrique
$x/2,y/2,+/2,-,+/2,y/2,x/2$. Une projection DG ne possède pas en général
d'inverse obtenu par un pas négatif. Cet ordre symétrique ne suffit donc
pas à établir l'ordre temporel deux à espaces fixés. Une diminution du
pas peut augmenter la diffusion de projection ou modifier les compensations
d'erreurs; la comparaison numérique inclut plusieurs nombres de pas.

\section{Observations et limites}
Sur la bosse fine du flot cellulaire à $t=0{,}005$, les supports $128^2 Q_4$
et quatre pas donnent une erreur $L^2$ d'environ $4{,}80\,10^{-4}$,
pour une meilleure projection FE à $4{,}79\,10^{-4}$.
Le support principal comporte $409600$ valeurs, contre $9025$ valeurs
pour le support Fourier initial; deux supports orientés de même dimension
interviennent temporairement. Ce résultat résout le défaut de précision
du contre-exemple, sans établir une supériorité générale de coût.

Le support $64^2 Q_4$ présente une erreur non monotone lorsque le nombre
de pas augmente. Une décomposition numérique en explique la compensation.
Soient $v_{2048}$ la solution à 2048 pas et $\Pi_Vu$ la projection FE de la
référence indépendante. Pour $a=v_{2048}-\Pi_Vu$ et $b=v_1-v_{2048}$,
les calculs GPU donnent
\[
 \|a\|^2=2{,}702\,10^{-7},\qquad
 \|b\|^2=2{,}876\,10^{-7},\qquad
 2\langle a,b\rangle=-5{,}147\,10^{-7}.
\]
Le terme croisé négatif réduit fortement l'erreur à un seul pas.
L'identité $\|a+b\|^2=\|a\|^2+\|b\|^2+2\langle a,b\rangle$ est
vérifiée à $5\,10^{-23}$ près. La solution à 2048 pas est un substitut
numérique du régime de petit pas, contrôlé par comparaison à 1024 pas,
et non une limite temporelle exacte. La bonne erreur à un seul pas
contient donc une compensation entre les deux écarts, sans établir
une amélioration de la consistance.
Le support $128^2 Q_4$ approche le plancher spatial entre
quatre et soixante-quatre pas. Les mesures de préparation incluent la
compilation lors de la première invocation. Les temps d'évolution sont
des mesures exploratoires et ne constituent pas une campagne comparative
avec intervalles de confiance.

La graduation utilise la position initiale connue. Le déplacement d'une
structure hors de la zone fine, le choix automatique des partitions,
leur adaptation en temps, les frontières et les vitesses FE générales
restent à étudier. Les projections conservatives, le transport SL-DG et
leurs réalisations GPU ont des antécédents établis, notamment
\href{https://arxiv.org/abs/2212.03036}{Einkemmer et Moriggl, méthode SL-DG GPU}
et les travaux de remappage conservatif cités dans le rapport principal.
La présente construction ne constitue pas, à elle seule, une revendication
de nouveauté pour un article.
\end{document}
